\documentclass[11pt]{article}
\usepackage{ap-physc}
\course{AP Physics C: Mechanics}
\headcourse{AP Physics C: Mechanics}
\headpart{Unit 3 \textperiodcentered\ Question Set}
\title{Unit 3: Work, Energy, and Power: Question Set}
\author{Paul Wang}
\date{November 2026}
\docpart{Unit 3 \textperiodcentered\ Question Set}

\begin{document}
\maketitle

\begin{abstract}
Exam-style practice on Unit 3 (topics 3.1--3.5), which is 15--25\% of the multiple-choice section. Section~1--4 hold 12 four-choice multiple-choice questions grouped by topic; most need a short calculation or a line of reasoning, as on the exam. The last page before the key is a full Qualitative/Quantitative Translation free-response question (8 points). All items are original and written for practice. Use the AP equation sheet and a calculator. Allow about 24 minutes for the multiple choice (the exam's pace is about 2 minutes per question) and 24 minutes for the free response. The answer key explains every item and names its topic; a scoring guide for the free response follows it.
\end{abstract}

\scoringline{MCQ: \blank\ / \qtotal \qquad FRQ: \blank\ / 8 \qquad Target: 8+ MCQ \qquad Pace: 2 min per MCQ}

\section{Work and kinetic energy (3.1--3.2)}

\begin{q}{1}{3.2 Dot product}{$W = \vec F\cdot\Delta\vec r = (2.0)(4.0) + (3.0)(-1.0) = 5.0$ J.}
A constant force $\vec F = (2.0\,\hat\imath + 3.0\,\hat\jmath)$ N acts on an object that moves through $\Delta\vec r = (4.0\,\hat\imath - 1.0\,\hat\jmath)$ m. How much work does the force do?
\chh{$5.0$ J}{$8.0$ J}{$11$ J}{$15$ J}
\end{q}
\begin{q}{2}{3.2 Work from F(x)}{$W = \int_1^3 5x^2\,dx = \tfrac53(27 - 1) = 43$ J.}
A force $F(x) = 5x^2$ (N, with $x$ in m) acts along the $x$-axis as an object moves from $x = 1.0$ m to $x = 3.0$ m. How much work does it do?
\chh{$40$ J}{$43$ J}{$50$ J}{$90$ J}
\end{q}
\begin{q}{3}{3.1 Work--energy theorem}{$\Delta K = W_{net} = 40$ J. $K_f = 9.0 + 40 = 49$ J $= \tfrac12(2.0)v^2$, so $v = 7.0$ m/s.}
A 2.0 kg block moving at 3.0 m/s is pushed by a net force of 10 N in its direction of motion for 4.0 m. What is its final speed?
\chh{$5.0$ m/s}{$6.3$ m/s}{$7.0$ m/s}{$9.3$ m/s}
\end{q}

\section{Potential energy (3.3)}

\begin{q}{4}{3.3 Force from potential energy}{$F = -dU/dx = -(6x - 6x^2)$. At $x = 2$: $-(12 - 24) = +12$ N.}
A particle's potential energy is $U(x) = 3x^2 - 2x^3$ (J, with $x$ in m). What is the force on the particle at $x = 2.0$ m?
\chh{$-12$ N}{$-4.0$ N}{$0$ N}{$+12$ N}
\end{q}
\begin{q}{3}{3.3 Equilibrium}{Equilibria are where $dU/dx = 6x(1 - x) = 0$: $x = 0$ and $x = 1$. Since $d^2U/dx^2 = 6 - 12x$ is $+6$ at $x = 0$ (a minimum), $x = 0$ is stable.}
For the same potential energy, $U(x) = 3x^2 - 2x^3$, at which position is the particle in stable equilibrium?
\chh{$x = 1.0$ m}{$x = 1.5$ m}{$x = 0$}{$x = 0.50$ m}
\end{q}
\begin{q}{1}{3.3 Gravitational potential energy}{$\Delta U = -\dfrac{GMm}{2R} - \left(-\dfrac{GMm}{R}\right) = \dfrac{GMm}{2R}$, and $GM = gR^2$, so $\Delta U = \tfrac12 mgR$.}
How much work is needed to lift a satellite of mass $m$ from Earth's surface (radius $R$) to a height $R$ above the surface? Express the answer in terms of $g$ at the surface.
\chh{$\tfrac12 mgR$}{$mgR$}{$2mgR$}{$\tfrac14 mgR$}
\end{q}

\section{Conservation of energy (3.4)}

\begin{q}{2}{3.4 Spring launch}{$\tfrac12 kx^2 = mgh$: $\tfrac12(200)(0.10)^2 = 1.0$ J, so $h = 1.0/(0.050 \times 9.8) = 2.0$ m.}
A spring ($k = 200$ N/m) compressed $0.10$ m launches a $0.050$ kg ball straight up. Ignoring air resistance, how high above its launch position does the ball rise? ($g = 9.8\ \mathrm{m/s^2}$)
\chh{$1.0$ m}{$2.0$ m}{$4.1$ m}{$41$ m}
\end{q}
\begin{q}{4}{3.4 Friction and energy}{All the gravitational energy becomes thermal energy on the floor: $mgh = \mu_k mg d$, so $d = h/\mu_k = 20$ m.}
A block slides from rest down a frictionless ramp from a height of 5.0 m, then across a level floor with $\mu_k = 0.25$. How far along the floor does it slide before stopping?
\chh{$1.3$ m}{$5.0$ m}{$10$ m}{$20$ m}
\end{q}
\begin{q}{1}{3.4 Energy and forces together}{Energy gives $v^2 = 2gL$ at the bottom. Then $T - mg = mv^2/L = 2mg$, so $T = 3mg$.}
A pendulum of length $L$ is released from rest with the string horizontal. What is the tension in the string when the bob passes through the lowest point?
\chh{$3mg$}{$mg$}{$2mg$}{$5mg$}
\end{q}
\begin{q}{2}{3.4 Escape speed}{Escaping means reaching $r \to \infty$ with $K \ge 0$: $\tfrac12 mv^2 - GMm/R = 0$.}
What is the escape speed from the surface of a planet of mass $M$ and radius $R$?
\chh{$\sqrt{GM/R}$}{$\sqrt{2GM/R}$}{$2GM/R$}{$\sqrt{GM/(2R)}$}
\end{q}

\section{Power (3.5)}

\begin{q}{3}{3.5 Power}{At constant speed the cable pulls with $mg = 7840$ N, so $P = Fv = 7840 \times 2.0 = 15.7$ kW.}
An elevator motor lifts an 800 kg car at a constant $2.0$ m/s. What power does the motor deliver? ($g = 9.8\ \mathrm{m/s^2}$)
\chh{$1.6$ kW}{$7.8$ kW}{$16$ kW}{$31$ kW}
\end{q}
\begin{q}{3}{3.5 Instantaneous power}{$v = (F/m)t = 15$ m/s at $t = 3.0$ s, so $P = Fv = 150$ W.}
A constant net force of 10 N acts on a 2.0 kg object that starts from rest. What power does the force deliver at $t = 3.0$ s?
\chh{$30$ W}{$75$ W}{$150$ W}{$225$ W}
\end{q}

\section{Free response}

\begin{frq}{Question 1 \textperiodcentered\ Qualitative/Quantitative Translation \textperiodcentered\ 8 points \textperiodcentered\ 24 minutes}
A block of mass $m$ is released from rest at height $h$ on a frictionless curved track. At the bottom the track becomes level, and the block slides into a horizontal spring of spring constant $k$ attached to a wall. The level section is also frictionless.

\begin{enumerate}[label=(\alph*),leftmargin=2em,itemsep=4pt]
\item The block is now released from height $2h$ instead. Without writing or deriving equations, explain how the spring's maximum compression changes. Say whether it doubles, more than doubles, or less than doubles, and why. \pts{3}
\item Derive an expression for the maximum compression $x_{max}$ in terms of $m$, $h$, $k$, and physical constants. \pts{3}
\item Explain how your expression in part (b) is consistent with your reasoning in part (a). \pts{2}
\end{enumerate}
\end{frq}

\makekey

\subsection*{Free-response scoring guide (8 points)}
\small
Award each point only when the answer does what the model does. Points are independent unless a part says otherwise; a correct method with an answer carried forward from an earlier error earns the later points.
\begin{scoring}
(a) (3) & 1 pt: doubling the height doubles the gravitational energy converted, so the spring must store twice the energy at maximum compression. 1 pt: the spring's stored energy grows faster than its compression (it depends on the compression squared, since each extra bit of compression takes a larger force). 1 pt: so the compression increases but by less than a factor of 2. & Saying the compression doubles because the energy doubles. \\
(b) (3) & 1 pt: energy conservation from release to maximum compression, with zero kinetic energy at both ends. 1 pt: $mgh = \tfrac12 kx_{max}^2$. 1 pt: $x_{max} = \sqrt{2mgh/k}$. & Using force balance, $kx = mg$, which describes a hanging spring, not this situation. \\
(c) (2) & 1 pt: $x_{max} \propto \sqrt h$, so doubling $h$ multiplies $x_{max}$ by $\sqrt2 \approx 1.41$. 1 pt: this is an increase of less than double, as argued in (a), because the stored energy $\tfrac12 kx^2$ depends on $x$ squared. & Restating the expression without connecting $\sqrt h$ to the claim in (a). \\
\end{scoring}

\end{document}
